\documentclass[11pt,a4paper]{article}
\usepackage[T1]{fontenc}
\usepackage[utf8]{inputenc}
\usepackage[ngerman]{babel}
\usepackage[a4paper,margin=2.5cm]{geometry}
\usepackage{setspace}
\usepackage{enumitem}
\usepackage{lmodern}
\usepackage[hidelinks]{hyperref}
\setlength{\parindent}{0pt}
\setlength{\parskip}{0.75em}
\linespread{1.04}
\setlist[itemize]{leftmargin=*,nosep}
\newcommand{\fillline}{\makebox[4.5cm]{\dotfill}}

\begin{document}
\setstretch{1.04}
\thispagestyle{empty}

\begin{center}
{\LARGE Vertrag zur Auftragsverarbeitung gemäß Art. 28 DSGVO\par}
\end{center}

\vspace{1em}

zwischen\par
\vspace{1.5em}
\textbf{................................................................................................}\par
\vspace{1em}
nachfolgend \glqq AN\grqq\ (Auftragnehmer) genannt\par
\vspace{1em}
und\par
\vspace{1em}
\textbf{................................................................................................}\par
\vspace{1em}
nachfolgend \glqq AG\grqq\ (Auftraggeber) genannt\par
\vspace{1em}
gemeinsam auch \glqq Parteien\grqq

\section*{1 Gegenstand und Dauer der Verarbeitung}
(1) Gegenstand dieser Vereinbarung ist die Verarbeitung personenbezogener Daten durch den AN im Auftrag des AG im Rahmen des zwischen den Parteien geschlossenen Hosting-, Pflege-, Support- und sonstigen Betreuungsvertrags bzw. der hierzu gehörenden technischen Leistungen vom \fillline .

(2) Die Auftragsverarbeitung umfasst insbesondere Hosting, Serverbetrieb, Bereitstellung von Webanwendungen, Protokollierung, Backup, Wiederherstellung, Wartung, Fernzugriff, Sicherheitsmaßnahmen sowie den Einsatz vorgeschalteter DNS-, Proxy-, SSL/TLS-, WAF-, DDoS-Schutz-, Cache-, CDN-, Rate-Limiting-, Bot-Schutz- und Turnstile-Dienste.

(3) Die Verarbeitung beginnt mit Aufnahme der Leistungen und endet mit deren Beendigung sowie der vollständigen Rückgabe oder Löschung der personenbezogenen Daten nach Maßgabe dieser Vereinbarung.

\section*{2 Art und Zweck der Verarbeitung}
(1) Der AN verarbeitet personenbezogene Daten ausschließlich zum Betrieb, zur Bereitstellung, Absicherung, Wartung, Fehleranalyse, Protokollierung, Datensicherung, Wiederherstellung, technischen Administration sowie zur vorgeschalteten Schutz- und Routingverarbeitung der vom AG betriebenen Website und der zugehörigen Systeme.

(2) Die Verarbeitung umfasst insbesondere IP-Adressen, Server- und Anwendungslogs, Formulareingaben, Datenbankinhalte, Admin-Zugriffe, Sicherungskopien, Wiederherstellungsvorgänge sowie die Verarbeitung über vorgeschaltete Schutz- und Routingdienste wie DNS, Reverse Proxy, SSL/TLS-Terminierung, Web Application Firewall, DDoS-Schutz, Cache, CDN, Rate Limiting, Bot-Schutz und Turnstile, soweit diese für die Leistungserbringung aktiviert sind.

(3) Eine Verarbeitung zu eigenen Zwecken des AN findet nicht statt, soweit nicht zwingende gesetzliche Pflichten entgegenstehen.

\section*{3 Art der personenbezogenen Daten und Kategorien betroffener Personen}
(1) Gegenstand der Verarbeitung sind je nach konkret vereinbartem Projekt und technischem Setup ausschließlich die für Hosting, Absicherung, Wartung und Support typischerweise erforderlichen personenbezogenen Daten. Hierzu gehören insbesondere IP-Adressen, Zeitstempel, Server-, Zugriffs- und Sicherheitsprotokolle, Kontakt- und Kommunikationsdaten aus Website-Formularen oder Supportvorgängen sowie Administrations- und Zugangsdaten.

(2) Besondere Kategorien personenbezogener Daten nach Art. 9 DSGVO sind nicht Gegenstand dieser Vereinbarung, es sei denn, ihre Verarbeitung ist im Hauptvertrag ausdrücklich vorgesehen oder wird vom AG vorab schriftlich oder in Textform dokumentiert angewiesen. In diesem Fall ist Voraussetzung, dass der AG die datenschutzrechtliche Zulässigkeit sowie die Erforderlichkeit erhöhter Schutzmaßnahmen vorab sichergestellt hat.

(3) Eine Verarbeitung weiterer personenbezogener Datenkategorien erfolgt nur nach vorheriger dokumentierter Weisung des AG und nach gesonderter Prüfung der technischen und organisatorischen Umsetzbarkeit durch den AN.

(4) Betroffene Personen sind insbesondere Website-Besucher, Kommunikations- und Kontaktpersonen, Nutzer von Website-Formularen sowie Mitarbeiter, Administratoren oder sonstige vom AG autorisierte Nutzer der vom AG betriebenen Systeme.

\section*{4 Weisungsrecht und Verantwortlichkeit}
(1) Der AN verarbeitet personenbezogene Daten ausschließlich auf dokumentierte Weisung des AG, sofern er nicht durch das Recht der Union oder der Mitgliedstaaten, dem der AN unterliegt, zu einer weitergehenden Verarbeitung verpflichtet ist.

(2) In einem Fall nach Abs. 1 teilt der AN dem AG diese rechtlichen Anforderungen vor der Verarbeitung mit, sofern das betreffende Recht eine solche Mitteilung nicht wegen eines wichtigen öffentlichen Interesses verbietet.

(3) Hält der AN eine Weisung des AG für rechtswidrig, weist er den AG unverzüglich darauf hin und ist berechtigt, die Ausführung bis zur Bestätigung oder Änderung der Weisung auszusetzen.

(4) Der AG ist für die Rechtmäßigkeit der Datenverarbeitung, die Zulässigkeit der Datenübermittlung an den AN sowie für die Inhalte seiner Website und die dort eingesetzten Formulare, Funktionen und Prozesse selbst verantwortlich.

\section*{5 Vertraulichkeit}
(1) Der AN stellt sicher, dass alle Personen, die auf personenbezogene Daten zugreifen können, zur Vertraulichkeit verpflichtet sind.

(2) Der Zugriff auf personenbezogene Daten ist auf das für die Vertragserfüllung erforderliche Maß zu beschränken.

(3) Der AN verpflichtet sich, personenbezogene Daten nur solchen Mitarbeitern oder Erfüllungsgehilfen zugänglich zu machen, die diese Daten zur Erfüllung ihrer Aufgaben benötigen.

\section*{6 Technische und organisatorische Maßnahmen}
(1) Der AN trifft die nach Art. 32 DSGVO erforderlichen technischen und organisatorischen Maßnahmen, um ein dem Risiko angemessenes Schutzniveau zu gewährleisten.

(2) Zu den Maßnahmen gehören insbesondere:
\begin{itemize}
  \item Zutritts-, Zugangs- und Zugriffskontrollen,
  \item rollenbasierte Berechtigungskonzepte,
  \item Mehrfaktor-Authentisierung für administrative Zugänge, soweit technisch verfügbar,
  \item verschlüsselte Übertragung personenbezogener Daten über aktuelle Transportverschlüsselung,
  \item Protokollierung administrativer Zugriffe,
  \item Maßnahmen zur Erkennung und Abwehr unbefugter Zugriffe,
  \item Backup- und Wiederherstellungsverfahren,
  \item Patch- und Schwachstellenmanagement,
  \item Maßnahmen zur Sicherstellung von Vertraulichkeit, Integrität, Verfügbarkeit und Belastbarkeit der Systeme.
\end{itemize}

(3) Der AN ist berechtigt, die technischen und organisatorischen Maßnahmen weiterzuentwickeln, soweit das vereinbarte Schutzniveau nicht unterschritten wird.

(4) Die jeweils konkret umgesetzten Maßnahmen ergeben sich ergänzend aus Anlage 1 zu dieser Vereinbarung. Der AN stellt sicher, dass die dort beschriebenen Maßnahmen im Rahmen des tatsächlich eingesetzten technischen Setups fortlaufend überprüft und bei Bedarf angepasst werden.

\section*{7 Unterauftragsverarbeiter}
(1) Der AG erteilt dem AN die allgemeine Genehmigung zum Einsatz von Unterauftragsverarbeitern nach Maßgabe dieser Vereinbarung.

(2) Der AN setzt die in Anlage 2 benannten Unterauftragsverarbeiter ein. Anlage 2 enthält die zum Zeitpunkt des Vertragsschlusses bekannten primären Unterauftragsverarbeiter sowie ergänzende Verweise auf die jeweils aktuellen, von diesen Anbietern veröffentlichten Listen weiterer Unterauftragsverarbeiter innerhalb der für die Leistungserbringung relevanten Verarbeitungskette.

(3) Zum Zeitpunkt des Vertragsschlusses sind dies insbesondere:
\begin{itemize}
  \item STRATO GmbH, Otto-Ostrowski-Straße 7, 10249 Berlin, Deutschland, für Server-, Hosting- und Infrastrukturleistungen;
  \item Cloudflare, Inc., 101 Townsend St., San Francisco, CA 94107, USA, für DNS-, Reverse-Proxy-, SSL/TLS-, WAF-, DDoS-Schutz-, Routing-, Cache-, CDN-, Rate-Limiting-, Bot-Schutz- und Turnstile-Leistungen, soweit diese aktiviert sind.
\end{itemize}

(4) Der AN hält dem AG die Identität der für die konkrete Leistungserbringung eingesetzten Unterauftragsverarbeiter in geeigneter Form aktuell verfügbar und informiert den AG auf Anfrage über die für das konkrete Setup maßgebliche Verarbeitungskette, soweit ihm diese Informationen vorliegen oder von den eingesetzten Anbietern veröffentlicht werden.

(5) Der AN schließt mit jedem von ihm unmittelbar eingesetzten Unterauftragsverarbeiter eine Vereinbarung, die diesem im Wesentlichen dieselben Datenschutzpflichten auferlegt, wie sie in dieser Vereinbarung zwischen AG und AN festgelegt sind.

(6) Der AN bleibt gegenüber dem AG für die Einhaltung der Datenschutzpflichten der eingesetzten Unterauftragsverarbeiter nach Maßgabe des Art. 28 Abs. 4 DSGVO verantwortlich.

(7) Über beabsichtigte Änderungen in Bezug auf von ihm unmittelbar eingesetzte Unterauftragsverarbeiter informiert der AN den AG mindestens 14 Tage vor deren Wirksamwerden in Textform.

(8) Soweit ein eingesetzter Anbieter seinerseits Änderungen in seiner Unterauftragsverarbeiterkette nach einem von ihm vorgegebenen Benachrichtigungs- oder Einwendungsprozess veröffentlicht, informiert der AN den AG hierüber unverzüglich nach Kenntnis oder durch Verweis auf die veröffentlichte Liste. Bei Cloudflare-bezogenen Änderungen hat der AG Einwendungen aus berechtigten datenschutzrechtlichen Gründen unverzüglich, spätestens jedoch innerhalb von 10 Tagen nach Zugang der Information des AN, mitzuteilen, damit der AN die vom Anbieter vorgesehene Einwendungsfrist wahren kann.

(9) Der AG kann aus berechtigten datenschutzrechtlichen Gründen innerhalb der jeweils anwendbaren Frist widersprechen. Widerspricht der AG und ist dem AN die Leistungserbringung ohne den betroffenen Unterauftragsverarbeiter nicht zumutbar möglich, sind die Parteien berechtigt, den Hauptvertrag hinsichtlich der betroffenen Leistungen außerordentlich zu kündigen.

\section*{8 Unterstützungspflichten und Meldungen}
(1) Der AN unterstützt den AG unter Berücksichtigung der Art der Verarbeitung und der ihm zur Verfügung stehenden Informationen bei der Erfüllung der Pflichten nach Art. 12 bis 22 sowie Art. 32 bis 36 DSGVO.

(2) Der AN leitet Anfragen betroffener Personen, die ihn unmittelbar erreichen, unverzüglich an den AG weiter und beantwortet diese nicht ohne Abstimmung mit dem AG.

(3) Der AN meldet dem AG Verletzungen des Schutzes personenbezogener Daten unverzüglich nach Bekanntwerden und stellt dem AG die zur Bewertung und Erfüllung gesetzlicher Meldepflichten erforderlichen Informationen zur Verfügung, insbesondere soweit verfügbar Art der Verletzung, betroffene Systeme und Datenkategorien, bereits ergriffene Gegenmaßnahmen sowie bekannte oder wahrscheinliche Auswirkungen.

(4) Der AN unterstützt den AG insbesondere bei Betroffenenrechten, Datenschutzverletzungen, Datenschutz-Folgenabschätzungen und etwaigen Konsultationen mit Aufsichtsbehörden.

(5) Soweit dem AN im Zusammenhang mit außergewöhnlichen, über die gewöhnliche Vertragserfüllung hinausgehenden Unterstützungshandlungen zusätzliche Kosten entstehen, kann eine gesonderte Vergütung vereinbart werden.

\section*{9 Löschung und Rückgabe}
(1) Nach Beendigung des Hauptvertrags löscht oder gibt der AN personenbezogene Daten nach Wahl des AG zurück, soweit keine gesetzlichen Aufbewahrungspflichten entgegenstehen.

(2) Produktivdaten werden nach entsprechender Weisung des AG gelöscht oder in einem gängigen elektronischen Format herausgegeben. Sicherungskopien, technische Zwischenspeicher, temporäre Exporte und sonstige redundante Kopien werden im Rahmen der regulären Backup-Retention oder nach Abschluss des jeweiligen technischen Zwecks gelöscht und danach nicht erneut produktiv genutzt.

(3) Auf Wunsch des AG bestätigt der AN die Löschung oder Rückgabe in Textform.

\section*{10 Kontrollrechte und Nachweise}
(1) Der AN weist dem AG auf Anfrage die Einhaltung dieser Vereinbarung in geeigneter Weise nach.

(2) Der AN kann den Nachweis der Einhaltung vorrangig durch aktuelle Zertifikate, Auditberichte, Prüfvermerke, TOM-Dokumentationen, Selbstauskünfte oder vergleichbare Nachweise erbringen.

(3) Weitergehende Vor-Ort-Kontrollen oder vergleichbare Prüfungen sind nur nach angemessener Vorankündigung und nur insoweit zulässig, als sie nach den vorgelegten Nachweisen erforderlich erscheinen und Betriebs- und Sicherheitsinteressen des AN sowie Rechte Dritter angemessen gewahrt werden.

(4) Kontrollen erfolgen unter Wahrung von Betriebs- und Geschäftsgeheimnissen sowie der Vertraulichkeit gegenüber Dritten.

\section*{11 Drittlandübermittlungen}
(1) Soweit im Rahmen der Leistungserbringung personenbezogene Daten außerhalb des Europäischen Wirtschaftsraums verarbeitet oder von dort aus zugänglich gemacht werden, stellt der AN sicher, dass die Voraussetzungen der Art. 44 ff. DSGVO eingehalten werden.

(2) Soweit Dienste von Cloudflare, Inc. eingesetzt werden, kann im Zusammenhang mit DNS-, Proxy-, Sicherheits-, Analyse-, Bot-Schutz-, Rate-Limiting-, Cache-, CDN- oder Turnstile-Funktionen ein Zugriff oder eine Verarbeitung personenbezogener Daten in Drittländern, insbesondere in den USA, nicht ausgeschlossen werden.

(3) Cloudflare wird, soweit Dienste zur Absicherung und Auslieferung der Website eingesetzt werden, je nach konkret aktiviertem Dienst als Unterauftragsverarbeiter des AN oder nach eigener Dokumentation in eigener datenschutzrechtlicher Verantwortlichkeit tätig. Soweit Cloudflare oder ein anderer eingesetzter Drittanbieter bestimmte Verarbeitungen in eigener datenschutzrechtlicher Verantwortlichkeit vornimmt, sind diese Verarbeitungen nicht Gegenstand dieser Vereinbarung; der AN schuldet insoweit weder eine Verantwortlichkeit als Auftragsverarbeiter noch eine vollständige Datenverarbeitung ausschließlich innerhalb der Europäischen Union, keine über die Anbieterunterlagen hinausgehenden Garantien und keine eigenständige rechtliche Prüfung des Drittanbieters.

(4) In diesem Fall stützt der AN die Übermittlung auf die jeweils einschlägigen zulässigen Transfermechanismen, insbesondere einen Angemessenheitsbeschluss, soweit anwendbar, die von dem Dienstanbieter vorgesehenen Standardvertragsklauseln sowie gegebenenfalls ergänzende Schutzmaßnahmen.

(5) Der AN informiert den AG auf Anfrage über die für die eingesetzten Dienste maßgeblichen Transfergrundlagen. Soweit bestimmte Regionalisierungs- oder Lokalisierungsfunktionen eines Dienstanbieters nicht vereinbart oder technisch nicht verfügbar sind, schuldet der AN keine ausschließliche Datenverarbeitung innerhalb der Europäischen Union.

\section*{12 Schriftform, Gerichtsstand, anwendbares Recht}
(1) Änderungen und Ergänzungen dieser Vereinbarung bedürfen der Textform.

(2) Es gilt deutsches Recht.

(3) Sofern der AG Kaufmann ist oder als juristische Person des öffentlichen Rechts bzw. als öffentlich-rechtliches Sondervermögen handelt, ist Gerichtsstand der Sitz des AN.

\section*{13 Schlussbestimmungen}
Sollten einzelne Bestimmungen dieser Vereinbarung ganz oder teilweise unwirksam sein oder werden, bleibt die Wirksamkeit der übrigen Bestimmungen unberührt.

\section*{Anlage 1: Technische und organisatorische Maßnahmen}
\textbf{1. Zutrittskontrolle}
\begin{itemize}
  \item Schutz der Räumlichkeiten und Systeme vor unbefugtem physischen Zugriff
  \item Zugriff auf administrative Endgeräte und Zugangsdaten nur für berechtigte Personen
  \item Sperrung oder Absicherung genutzter Endgeräte bei Abwesenheit
\end{itemize}

\textbf{2. Zugangskontrolle}
\begin{itemize}
  \item Personalisierte Benutzerkonten
  \item Starke Passwörter
  \item Mehrfaktor-Authentisierung für administrative Zugänge, soweit technisch verfügbar und im eingesetzten Dienst aktivierbar
  \item Regelmäßige Überprüfung aktiver Zugänge
  \item Nutzung verschlüsselter Verbindungen für administrative Logins
\end{itemize}

\textbf{3. Zugriffskontrolle}
\begin{itemize}
  \item Rollen- und Berechtigungskonzepte nach dem Need-to-know-Prinzip
  \item Beschränkung administrativer Rechte auf den erforderlichen Umfang
  \item Protokollierung administrativer Zugriffe und Änderungen
  \item Nutzung getrennter Zugänge für unterschiedliche Dienste und Kundenkontexte, soweit technisch möglich
\end{itemize}

\textbf{4. Weitergabekontrolle}
\begin{itemize}
  \item Verschlüsselte Übertragung personenbezogener Daten über aktuelle Transportverschlüsselung
  \item Beschränkung externer Datenzugriffe auf autorisierte Systeme und Personen
  \item Prüfung eingesetzter Dienstleister und Schnittstellen
  \item Dokumentierte Freigabe oder Beauftragung externer Übermittlungen, soweit erforderlich
\end{itemize}

\textbf{5. Eingabekontrolle}
\begin{itemize}
  \item Nachvollziehbarkeit administrativer Änderungen
  \item Protokollierung sicherheitsrelevanter Eingriffe, soweit technisch möglich
  \item Dokumentation wesentlicher Konfigurationsänderungen und Wiederherstellungsvorgänge
\end{itemize}

\textbf{6. Verfügbarkeitskontrolle}
\begin{itemize}
  \item Regelmäßige Backups oder Snapshots nach betrieblicher Erforderlichkeit
  \item Dokumentierte Wiederherstellungsverfahren
  \item Schutz vor Datenverlust durch technische Defekte, Fehlbedienung oder Angriffe
  \item Durchführung von Wiederherstellungsmaßnahmen im Störungsfall nach Priorität und technischer Verfügbarkeit
\end{itemize}

\textbf{7. Trennungskontrolle}
\begin{itemize}
  \item Logische Trennung von Umgebungen und Kundenkontexten, soweit technisch erforderlich
  \item Trennung von Produktiv-, Test- und Administrationszugängen, soweit vorhanden
  \item Verwendung getrennter Datenbestände oder Zugriffsbereiche für Test- und Produktivumgebungen, soweit technisch umsetzbar
\end{itemize}

\textbf{8. Auftragskontrolle}
\begin{itemize}
  \item Verarbeitung nur auf dokumentierte Weisung
  \item Vertragliche Bindung eingesetzter Unterauftragsverarbeiter
  \item Regelmäßige Überprüfung der datenschutzrelevanten Prozesse
  \item Dokumentation der eingesetzten Unterauftragsverarbeiter und ihrer Funktionen
\end{itemize}

\textbf{9. Schutz vor unbefugtem Zugriff Dritter}
\begin{itemize}
  \item Firewall- und Schutzmechanismen gegen Angriffe
  \item Patch- und Schwachstellenmanagement
  \item Maßnahmen zur Erkennung und Abwehr missbräuchlicher Zugriffe
  \item Absicherung administrativer Zugänge, insbesondere über verschlüsselte Protokolle und beschränkte Berechtigungen
\end{itemize}

\textbf{10. Integrität und Vertraulichkeit der Endgeräte}
\begin{itemize}
  \item Aktuelle System- und Sicherheitsupdates auf administrativ genutzten Endgeräten
  \item Einsatz von Geräte- oder Datenträgerverschlüsselung, soweit technisch verfügbar
  \item Schutz lokaler Zugangsdaten und Konfigurationsdateien vor unbefugtem Zugriff
\end{itemize}

\textbf{11. Dokumentation und Überprüfung}
\begin{itemize}
  \item Regelmäßige Überprüfung und Aktualisierung der Schutzmaßnahmen
  \item Anpassung der Maßnahmen bei technischen oder organisatorischen Änderungen
  \item Überprüfung, ob die tatsächliche Verarbeitung weiterhin durch die vereinbarten TOMs abgedeckt ist
\end{itemize}

\section*{Anlage 2: Unterauftragsverarbeiter}
\textbf{1. STRATO GmbH, Otto-Ostrowski-Straße 7, 10249 Berlin, Deutschland}\par
Leistung: Server-, Hosting- und Infrastrukturleistungen\par
Ergänzender Verweis: Weitere von STRATO eingesetzte Auftragsverarbeiter ergeben sich, soweit für das konkrete Produkt oder Setup relevant, aus der jeweils aktuellen \href{https://www.strato.de/agb/avv/}{STRATO-AVV samt Anhang}.

\textbf{2. Cloudflare, Inc., 101 Townsend St., San Francisco, CA 94107, USA}\par
Leistung: DNS-, Reverse-Proxy-, SSL/TLS-, WAF-, DDoS-Schutz-, Routing-, Cache-, CDN-, Rate-Limiting-, Bot-Schutz- und Turnstile-Leistungen, soweit aktiviert\par
Ergänzender Verweis: Weitere von Cloudflare eingesetzte Unterauftragsverarbeiter ergeben sich, soweit für die konkret aktivierten Dienste relevant, aus der jeweils aktuellen \href{https://www.cloudflare.com/gdpr/subprocessors/cloudflare-services/}{Cloudflare-Sub\-pro\-zes\-sor\-liste}.\par
Hinweis: Bei Nutzung der Dienste kann eine Verarbeitung personenbezogener Daten in Drittländern, insbesondere in den USA, nicht ausgeschlossen werden. Eine ausschließliche Verarbeitung innerhalb der EU wird nur geschuldet, soweit entsprechende Regionalisierungsfunktionen ausdrücklich vereinbart und tatsächlich eingesetzt werden.\par
Zusatzhinweis zu Turnstile: Soweit Cloudflare Turnstile eingesetzt wird, ist zu beachten, dass Cloudflare nach eigener Dokumentation bestimmte Signale nicht nur als Auftragsverarbeiter zur Absicherung der Website, sondern zusätzlich in eigener datenschutzrechtlicher Verantwortlichkeit zur Verbesserung der Bot-Erkennung verarbeitet. Diese eigenverantwortliche Verarbeitung von Cloudflare ist nicht Gegenstand dieser AVV zwischen AG und AN.

\vfill
\begin{tabular}{p{0.45\textwidth}p{0.45\textwidth}}
\hrulefill & \hrulefill \\
Ort, Datum & Ort, Datum \\
\\[3cm]
\hrulefill & \hrulefill \\
Unterschrift AN & Unterschrift AG \\
\end{tabular}

\end{document}
